\clearpage\thispagestyle{empty}
\pdfbookmark[0]{Quellen zur Rezension der Reise von Humboldt und Bonpland}{voyage1809-sources}
\begin{center}\textbf{Quellen zur Rezension der Reise von Humboldt und Bonpland}\end{center}
\noindent Diese KI-übersetzte Fassung in modernem Deutsch folgt der deutschsprachigen Rezension in Carl Friedrich Gauß, \emph{Werke}, Band~VI (Göttingen, 1874), S.~512–515. Sie erschien zuerst ohne Verfassernamen unter Paris und Tübingen in den \emph{Göttingischen gelehrten Anzeigen}, Stück~91, 10. Juni 1809, S.~898–904. Die Niedersächsische Staats- und Universitätsbibliothek Göttingen stellt die \href{https://images.p.sub.uni-goettingen.de/iiif/image/gdz:PPN319721507_1809:00000900/full/1750,/0/default.jpg}{Anfangsseite}, die \href{https://images.p.sub.uni-goettingen.de/iiif/image/gdz:PPN319721507_1809:00000903/full/1750,/0/default.jpg}{Koordinatentabelle} und die \href{https://images.p.sub.uni-goettingen.de/iiif/image/gdz:PPN319721507_1809:00000906/full/1750,/0/default.jpg}{Schlussseite} dieses Drucks bereit. Der Erstdruck wiederholt die Einheiten in beiden Zeilen und verwendet Striche für Zeitminuten und Zeitsekunden. Die Werkausgabe gibt die Einheiten in der ersten Zeile an und lässt sie in der zweiten unausgeschrieben. Im Erstdruck steht \emph{die} bzw. \emph{selbst}, in der Werkausgabe \emph{die die} bzw. \emph{selbs}.
\par\smallskip\noindent Der Buchtitel bezeichnet den vierten Teil der \emph{Reise von Humboldt und Bonpland}, \emph{Astronomie und Magnetismus}. Der erste Band enthält eine Sammlung astronomischer Beobachtungen, trigonometrischer Operationen und barometrischer Messungen, die auf einer Reise in die Äquatorialregionen des Neuen Kontinents von 1799 bis 1803 vorgenommen wurden. Die Rezension behandelt die beiden Lieferungen von 1808 und 1809, die bei Schöll und Cotta erhältlich waren. Alexander von Humboldt stellte die Beobachtungen an, Aimé Bonpland begleitete ihn und Jabbo Oltmanns arbeitete die Ergebnisse aus. Der \href{https://www.digitale-sammlungen.de/view/bsb10868686?page=9}{gebundene erste Band in der Bayerischen Staatsbibliothek} trägt das spätere Impressum Paris, F.~Schoell, 1810. Humboldts \emph{Abhandlung über die Refraktionen in der heißen Zone} erschien auch im \href{https://annales.minesparis.psl.eu/articles/1808-2/89-114.pdf}{\emph{Journal des mines}, Band~24 (1808), S.~169–218}. Satz mit KI erstellt.
\clearpage
